Define a germ

Germs are an equivalence relation with respect to some restriction function.
For example a germ of functions \( X \to Y \), are all functions which coincide, when restricted 
to some subset \( U \) of \( X \). Typically one takes a point \( x \in X \), and says
that \( f \sim g \) (are germ equivalent), iff. \( \exists U, x \in U, f(u) = g(u), \forall u \in U \).

In particular a local properties would be shared by such functions, such as being continuous or analytic.

In complex analysis, the subject of analytic continuation deals with germs.
Namely let \( f \) be some analytic function in the disc \( D_r(z_0) \), so 
\( f(z)=\sum_{k=0}^\infty \alpha_k (z-z_0)^k \)
converges. We call the vector \( g = (z_0, \alpha_0, \alpha_1, \alpha_2, \ldots) \) a germ.


Define a presheafs and sheafs 

Let \(X\) be a topological space. A presheaf \(\mathcal{F}\) of sets on \(X\) consists of the following data:
1.For each open set \(U\subseteq X\), there exists a set \(\mathcal{F}(U)\). This set is also denoted \(\Gamma(U, \mathcal{F})\). 
  The elements in this set are called the sections of \(\mathcal{F}\) over \(U\). 
  The sections of \(\mathcal{F}\) over \(X\) are called the global sections of \(\mathcal{F}\).

2. For each inclusion of open sets \(V \subseteq U\), a function \(\operatorname{res}^{U}_{V} \colon \mathcal{F}(U) \rightarrow \mathcal{F}(V)\) exists. 
   In view of many of the examples below, the morphisms \(\text{res}^{U}_{V}\) are called restriction morphisms. 
   If \(s \in \mathcal{F}(U)\), then its restriction \(\text{res}^{U}_{V}(s)\) is often denoted \(s|_V\) by analogy with restriction of functions.
   The restriction morphisms are required to satisfy two additional functorial properties:

For every open set \(U\) of \(X\), the restriction morphism 
\(\operatorname{res}^{U}_{U} \colon \mathcal{F}(U) \rightarrow \mathcal{F}(U)\)
is the identity morphism on \(\mathcal{F}(U)\).
If we have three open sets \(W \subseteq V \subseteq U\), then the function composition \(\text{res}^{V}_{W}\circ\text{res}^{U}_{V}=\text{res}^{U}_{W}\).

Axiomatically, a sheaf is a presheaf that satisfies both of the following axioms:
1. (Locality) Suppose \(U\) is an open set, \(\{ U_i \}_{i \in I}\) is an open cover of \(U\) with \( U_i \subseteq U\) for all \( i \in I\), and \(s, t \in \mathcal{F}(U)\) 
    are sections. If \(s|_{ U_i} = t|_{ U_i}\) for all \(i \in I\), then \(s = t\).
2. (Gluing axiom) Suppose \(U\) is an open set, \(\{ U_i \}_{i \in I}\) is an open cover of \(U\) with \( U_i \subseteq U\) for all \( i \in I\), and 
    \(\{ s_i \in \mathcal{F}(U_i) \}_{i \in I}\) is a family of sections. If all pairs of sections agree on the overlap of their domains, that is, 
    if \(s_i|_{U_i\cap U_j} = s_j|_{U_i \cap U_j}\) for all \(i, j \in I\), then there exists a section \(s \in \mathcal{F}(U)\) such that \(s|_{U_i} = s_i\) for all \(i \in I\).

An example of a presheaf, which is not a sheaf: Let \( X = \mathbb{R} \) and \( \mathcal{F}(U) \) all cont. function which are constant on \( U \).
An example of a sheaf: Let \( X = \mathbb{R} \) and \( \mathcal{F}(U) \) all cont. function on \( U \).
Another example of a sheaf: A ringed space \( (X, \mathcal{O}_X) \) is a top. space \( X \) with a sheaf of rings \( \mathcal{O}_X \) on \( X \).
The sheaf is called the structure sheaf.
A locally ringed space is ringed space s.t. all stalks of \( \mathcal{O}_X \) are local rings.

More generally 

is defined to 

generalization of
